\section*{ID9102: Definition of ECas(L)}
\paragraph{Metadata.}
PiID: 0941130315095 | RTEID: RTE:P35345.0873:T1.6892 | UUID: 141975eb-4b42-5383-b6a9-9ad57be0c4c6

\paragraph{Canonical Statement.}
uences: Casimir pressures; boundaryinduced particle creation; modifications to correlation functions; interplay with vacuum polarization. Falsifiability / Test Handle: Experimental measurement of boundaryinduced forces or shifts that disagree with this operator’s predictions would test its validity. 0.0.58 ID 056: CasimirType Phenomena Type: Conservation Law Formal Statement: The Casimir energy between two uncharged, parallel boundaries separated by distance LLL in the Trawinistic framework is \$ECas(L)=π2720L3(1+)E \{\textbackslash{}text\{Cas\}\}(L)=\$ -\textbackslash{}frac\{\textbackslash{}piˆ2\}\{720 Lˆ3\} \$(1+\textbackslash{}epsilon)ECas(L)=720L3π2(1+),\$ where \textbackslash{}epsilon encodes curvature and TZP corrections. Interpretive Meaning: Even empty space exerts forces due to vacuum fluctuations; these forces are modified by the TZP environment. Dependencies: ID 055, ID 051 Derivable Consequences: Attraction between plates in vacuum; constraints on boundary configurations; contributions to renormal

\paragraph{Canonical Equation.}
\[
ECas(L)=π2720L3(1+)E {\text{Cas}}(L)=
\]

\paragraph{Dictionary.}
Definition of ECas(L) — ECas(L)=π2720L3(1+)E Cas(L)=

\paragraph{Encyclopedia.}
Definition of ECas(L) is an algebraic definition, written ECas(L)=π2720L3(1+)E Cas(L)=. It is an algebraic definition expressing the quantity directly in terms of the variables on its right-hand side. It is built from L, E, defining ECas(L). Within the Trawin Zero Point registry it serves as one element of the framework's mathematical narrative.
